% Converted from the authoritative recovery source, 2026-11-14.
\subsection{The original AA potential works without changing its
logic}\label{influence-recovery:the-original-aa-potential-works-without-changing-its-logic}

Under an arbitrary product of blocks, write the orthogonal Hoeffding
decomposition \(g=\sum_Sg_S\), where each \(g_S\) depends on \(S\) and
has mean zero in each block in \(S\). Global block degree at most
\(\Delta\) implies \(g_S=0\) for \(|S|>\Delta\). Therefore

\[
 \mathcal P(g):=\sum_a J_a(g)
 =\sum_S |S|\|g_S\|_2^2
 \le\Delta\operatorname{Var}(g)\le\Delta/4.
 \tag{P1}
\]

Revealing block \(a\) and continuing under its conditional product law
gives the exact identity

\[
 \mathbb E_{Z_a}\mathcal P(g\mid Z_a)
 =\mathcal P(g)-J_a(g).
 \tag{P2}
\]

For example, every component containing \(a\) loses one from its support
size; components become orthogonal after averaging over the revealed
block. Thus parity causes no problem: a block can reduce this potential
even when the immediately revealed block explains zero variance of \(g\)
itself.

At each nonconstant decision-tree leaf, query a block with largest
\(J_a\) for the restricted function; constant leaves stop. This uses no
information other than the previously revealed source blocks. Let
\(V_t\) be the expected conditional variance of \(g\) given the source
transcript through depth \(t\), and let \(P_t\) be the average potential
there. Restricted functions retain block degree at most \(\Delta\) and
remain bounded on the full remaining cube, so (B2) applies at every
leaf. For

\[
 h(v)=\frac{v^{5/2}}{[k+\log(\Delta/v)]^6},\qquad h(0)=0,
\]

the second derivative is positive on \((0,1/4]\): more generally
\(v^aL(v)^{-b}\), where \(L(v)=k+\log(\Delta/v)\), has second derivative

\[
 v^{a-2}L^{-b}\left[a(a-1)
       +\frac{b(2a-1)}L+\frac{b(b+1)}{L^2}\right]>0
\]

when \(a>1\) and \(b>0\). Equations (B2), (P2), and Jensen consequently
give

\[
 P_t-P_{t+1}\ge\frac c{\Delta^6}h(V_t).
 \tag{P3}
\]

\(V_t\) is nonincreasing. If \(V_s>e^2\), every one of the \(s\) steps
decreases the potential by at least
\(ce^5/[\Delta^6(k+\log(\Delta/e^2))^6]\). The initial potential is at
most \(\Delta/4\). Hence a deterministic depth bound

\[
 S(\Delta,e;k)
 =\left\lceil
 C_2\Delta^7[k+\log(\Delta/e^2)]^6 e^{-5}
 \right\rceil+1
 \tag{P4}
\]

forces \(V_s\le e^2\). Outputting the conditional mean yields mean
absolute error at most \(e\). Here \(0<e<1/2\); larger target errors
follow by using \(e=1/2\). There is no expected-depth truncation or
extra error factor in (P4).

This is the usual AA proof: influential source block, exact drop in
total source influence, repeated restrictions, and conditional-mean
output. Dyadic encodings were used only to prove (B2); the resulting
source tree queries the original blocks with their exact laws.
